\section{分布式语义：大模型的理论基础}

\subsection{从稀疏到稠密}

大语言模型的核心原理是\textbf{分布式语义}——用一个词周围的上下文来定义这个词的意义（Firth的经典观点）。

\begin{knowledgebox}{关键技术突破}
将语言从高维空间中的\textbf{稀疏表示}转化为\textbf{稠密的向量空间几何结构}，这是一个重大进步。它使得语言处理问题完全变成了\textbf{数学计算问题}——原来稀疏的共现空间不可计算，变成稠密向量空间后就可以用数学工具处理。
\end{knowledgebox}

\subsection{可证明的极限性质}

可以证明：只要数据量足够多、上下文足够长，这个向量空间就会趋近于\textbf{语义关系空间}。这解释了为什么当前研究都在追求更长的上下文和更多的数据。

\begin{importantbox}{语义趋近的含义}
一旦向量空间趋近了语言的语义关系，在一定意义下就可以做到``理解''，甚至可以做到\textbf{自反性}——对自己的思考进行思考。这在大语言模型中已经被观察到了。
\end{importantbox}

\subsection{本章小结}

大语言模型基于分布式语义原理，通过稀疏到稠密的向量空间转换，实现了语言的可计算化。但这种语义定义本身是近似的。


